\subsection{Proof of Part (1) of Theorem \ref{theorem}}\label{part1} 
%
The proof of this case is carried out by a different way from the other cases. 
We only need the asymptotic behavior of the magnitude function at small scale. 

Assume $n=3$. Put $a=d_{12}, b=d_{13}$ and $c=d_{23}$. 
%
\begin{lemma}\label{lemma_Mu_Md}
For $\la=1,2$ and $3$, $M_\la=\lim_{t\to 0^+}\frac{d^\la}{dt^\la}M_X(t)$ are given by 
\[
\begin{array}{rcl}
M_1&=&\displaystyle \frac{2abc}{-a^2-b^2-c^2+2ab+2bc+2ca}, \\[6mm]
%
M_2&=&\displaystyle \frac{2abc(b+c-a)(c+a-b)(a+b-c)}{{(-a^2-b^2-c^2+2ab+2bc+2ca)}^2}, \\[6mm]
%
M_3&=&\displaystyle \frac{abc\,P(a,b,c)}{{(-a^2-b^2-c^2+2ab+2bc+2ca)}^3},
\end{array}
\]
where 
\[\begin{array}{rcl}
P(a,b,c)&=&\displaystyle a^6+\dots+a^5b+\dots-13a^4b^2-\dots+9a^4bc+\dots+22a^3b^3+\dots \\[2mm]
%
&&\displaystyle -10a^3b^2c-\dots+30a^2b^2c^2.
\end{array}
\]
\end{lemma}

Since the denominators of $M_\la$ are $\de_{2}^{\la}$ they are positive by Proposition \ref{pos_Mpd}. 

\begin{proof}
If we put\footnote{$\si_\mu$'s make the expressions simpler than the elementary symmetric polynomials of $a,b$ and $c$ do. } $\sigma_\mu=a^\mu+b^\mu+c^\mu$ $(\mu\in\NN)$, we have 
%
\[\begin{array}{c}
\nu_0=\nu_1=\de_0=\de_1=0,\> \nu_2=\de_2=\si_1^{\,2}-2\si_2,\\[3mm]
%
\displaystyle \nu_3=-\si_1\si_2+2\si_3, \, \nu_4=\frac{\si_1\si_3}3+\frac{\si_2^{\,2}}4-\frac{7\si_4}6, \, \nu_5=-\frac{\si_1\si_4}{12}-\frac{\si_2\si_3}6+\frac{\si_5}2,\\[3mm]
%
\displaystyle \nu_3-\de_3=2abc, \> \nu_4-\de_4=-abc\,\si_1, \> \nu_5-\de_5=\frac{abc}4\left(\si_1^{\,2}+\frac{\si_2}3\right).
\end{array}\]
%
Now $M_1$ is obtained from \eqref{M1} and $M_2$ from 
\[
M_2=2\,\frac{(\nu_4-\de_4)\de_2-(\nu_3-\de_3)\de_3}{\de_{2}^{\,2}}.
\]
$M_3$ is obtained in the same way, although the calculation is more complicated, so we omit the details. 
\end{proof}



\begin{lemma}\label{abc<-M1M2M3}
The edge lengths $a,b$ and $c$ can be obtained from $M_1,M_2$ and $M_3$. 
\end{lemma}

\begin{proof}
Let $s_1, s_2$ and $s_3$ be elementary symmetric polynomials of 
\begin{equation}\label{xyz<->abc}
x=b+c-a, \> y=c+a-b, \> z=a+b-c \,; \notag
\end{equation}
%
\begin{equation}\label{s123<->xyz}
\left\{\begin{array}{rcl}
s_1&=&x+y+z,\\[1mm]
s_2&=&xy+yz+zx, \\[1mm]
s_3&=&xyz.
\end{array}\right. \notag
\end{equation}
%
Since $x,y,z\ge0$ and at most one of $x,y$ and $z$ can be $0$, $s_1,s_2>0$ and $s_3\ge0$. 
%

Since $M_1,M_2$ and $M_3$ are symmetric in $a,b$ and $c$, and the elementary symmetric polynomials of $a, b$ and $c$ can be expressed by $s_1, s_2$ and $s_3$ as 
\begin{equation}\label{abcs1s2s3}
\left\{\begin{array}{rcl}
a+b+c&=&s_1, \\[2mm]
ab+bc+ca&=&\displaystyle \frac{s_1^{\,2}+s_2}4, \\[4mm]
abc&=&\displaystyle \frac{s_1s_2-s_3}8,
\end{array}\right.
\end{equation}
%
$M_1,M_2$ and $M_3$ can be expressed by $s_1, s_2$ and $s_3$; 
\[%
\begin{array}{rcl}
M_1&=&\displaystyle \frac{s_1s_2-s_3}{4s_2},\\[4mm]
%
M_2&=&\displaystyle \frac{(s_1s_2-s_3)s_3}{4s_2^{\,2}},\\[4mm]
%
M_3&=&\displaystyle -\frac{(s_1s_2-s_3)\left(s_1^{\,2}s_2^{\,2}+4s_1s_2s_3-3s_2^{\,3}-12s_3^{\,2}\right)}{32s_2^{\,3}}.
\end{array}%
\]
%
Remark that $M_1>0$ since $s_1s_2-s_3=8abc>0$ and $s_2>0$. 
Solving the above equations for $s_1, s_2$ and $s_3$, we obtain 
%
\begin{equation}\label{s123<->M123}
\begin{array}{rcl}
s_1&=&\displaystyle \frac{4M_1^{\,2}+M_2}{M_1}, \\[6mm]
%
s_2&=&\displaystyle \frac{16M_1^{\,4}+24M_1^{\,2}M_2+8M_1M_3-7M_2^{\,2}}{3M_1^{\,2}}, \\[6mm]
%
s_3&=&\displaystyle \frac{M_2\left(16M_1^{\,4}+24M_1^{\,2}M_2+8M_1M_3-7M_2^{\,2}\right)}{3M_1^{\,3}}.
\end{array}
\end{equation}
%
The edge lengths $a,b$ and $c$ are obtained from $s_1, s_2$ and $s_3$ by \eqref{abcs1s2s3}, and hence from $M_1,M_2$ and $M_3$ by \eqref{s123<->M123}, which completes the proof. 
\end{proof}

\begin{corollary}
A $3$-point metric space $X$ is determined by  $\displaystyle \lim_{t\to 0^+}M^{(\la)}_X(t)$ $(\la=1,2,3)$. 
\end{corollary}
